% =================================================================
% Abschnitt "Punktschätzer für den Erwartungswert: Rückblick"
% ===================================================================

\newglossaryentry{bem_poisson_parameter_aus_ew}{
    name={Wie lässt sich bei Poisson-verteilten Stichprobenvariablen der Parameter $\lambda$ aus dem Erwartungswert gewinnen?},
    description={$\lambda = E[X_1]$ -- der Parameter ist hier selbst der Erwartungswert.},
    type=dinge
}

\newglossaryentry{bem_exponential_parameter_aus_ew}{
    name={Wie lässt sich bei exponentialverteilten Stichprobenvariablen der Parameter $\alpha$ aus dem Erwartungswert gewinnen?},
    description={$\alpha = \dfrac{1}{E[X_1]}$ -- der Parameter ist der Kehrwert des Erwartungswerts.},
    type=dinge
}

\newglossaryentry{bem_parameter_nicht_immer_aus_ew_gewinnbar}{
    name={Lässt sich der Parameter einer Verteilung immer aus dem Erwartungswert berechnen?},
    description={Nein. Beispiel: Sind die Stichprobenvariablen gleichverteilt auf $\{-\theta,\dots,\theta\}$, so gilt stets $E[X_i]=0$ -- unabhängig von $\theta$. Der Erwartungswert enthält hier gar keine Information über $\theta$.},
    type=dinge
}

\newglossaryentry{bem_alternative_wenn_ew_nicht_ausreicht}{
    name={Was kann man tun, wenn der Erwartungswert nichts über den gesuchten Parameter verrät (wie bei der symmetrischen Gleichverteilung)?},
    description={Man kann stattdessen die Varianz schätzen, die (anders als der Erwartungswert) vom gesuchten Parameter abhängt.},
    type=dinge
}

% ===================================================================
% Abschnitt "Wiederholung: Schätzer für die Varianz"
% ===================================================================

\newglossaryentry{satz_varianzschaetzer_bekannter_ew}{
    name={Wie lautet der erwartungstreue Schätzer für die Varianz, wenn der Erwartungswert $\mu$ bekannt ist?},
    description={$\dfrac1n\sum_{i=1}^n (x_i-\mu)^2$ ist erwartungstreu -- man ersetzt $\bar x$ im naiven Ansatz durch den bekannten wahren Erwartungswert $\mu$.},
    type=dinge
}

% ===================================================================
% Abschnitt "Übergang: die Normalverteilung"
% ===================================================================

\newglossaryentry{bem_warum_normalverteilung_eingefuehrt}{
    name={Warum wird an dieser Stelle die Normalverteilung als Beispiel eingeführt?},
    description={Um zu verstehen, wie ein Schätzer für eine Varianz sinnvoll sein kann. Die Normalverteilung unterscheidet sich von den bisherigen Beispielen dadurch, dass sie durch zwei Parameter ($\mu$ und $\sigma^2$) beschrieben wird.},
    type=dinge
}

\newglossaryentry{def_dichte_normalverteilung_allgemein}{
    name={Wie lautet die Dichte der Normalverteilung mit Parametern $\mu\in\mathbb R$, $\sigma^2>0$?},
    description={$\varphi_{\mu,\sigma^2}(x) = \dfrac{1}{\sqrt{2\pi\sigma^2}} \exp\!\left(-\dfrac{(x-\mu)^2}{2\sigma^2}\right)$ für $x\in\mathbb R$. (Für $\mu=0,\sigma^2=1$ ergibt sich die Dichte der Standardnormalverteilung.)},
    type=dinge
}

\newglossaryentry{bem_fakten_ew_var_normalverteilung}{
    name={Welche Kennzahlen einer Zufallsvariablen $X$ mit Dichte $\varphi_{\mu,\sigma^2}$ entsprechen den beiden Parametern $\mu$ und $\sigma^2$?},
    description={$E[X]=\mu$ und $\operatorname{Var}(X)=\sigma^2$ -- die beiden Parameter der Normalverteilung sind also direkt Erwartungswert und Varianz.},
    type=dinge
}

% ===================================================================
% Abschnitt "Bereichsschätzer: Konfidenzintervalle"
% ===================================================================

\newglossaryentry{bem_motivation_bereichsschaetzer}{
    name={Was motiviert die Einführung von Bereichsschätzern (Konfidenzintervallen)?},
    description={Punktschätzungen (auch simulativ illustriert) zeigen, dass der geschätzte Parameter in der Regel nicht exakt der wahre Parameter ist. Die Idee ist daher, ein Intervall zu suchen, von dem man mit vorgegebener Wahrscheinlichkeit sagen kann, dass der wahre Parameter darin liegt.},
    type=dinge
}

\newglossaryentry{def_quantil_standardnormalverteilung_ki}{
    name={Wie ist das $(1-\alpha/2)$-Quantil $z_{1-\alpha/2}$ der Standardnormalverteilung definiert?},
    description={Das kleinste $c$ mit $\Phi(c)\ge 1-\alpha/2$ -- äquivalent dazu, dass für $Z\sim N(0,1)$ gilt $P(|Z|\le c)\ge 1-\alpha$. Es liefert die Breite des symmetrischen Bereichs, der eine standardnormalverteilte ZV mit Wahrscheinlichkeit $\ge 1-\alpha$ enthält.},
    type=dinge
}

\newglossaryentry{bem_darstellung_normalverteilte_zv_via_standardnormal}{
    name={Wie lässt sich eine normalverteilte Zufallsvariable $X\sim N(\mu,\sigma^2)$ mit Hilfe einer standardnormalverteilten ZV $Z$ darstellen?},
    description={$X = \mu + \sigma Z$ mit $Z\sim N(0,1)$.},
    type=dinge
}

\newglossaryentry{satz_summe_unabhaengiger_normalverteilter_zv_ki}{
    name={Welcher Fakt über Summen unabhängiger normalverteilter Zufallsvariablen wird zur Herleitung des Konfidenzintervalls für $\bar X$ genutzt?},
    description={Die Summe unabhängiger normalverteilter ZV ist wieder normalverteilt, wobei sich die Parameter addieren: $\sum_{i=1}^n X_i \sim N(n\mu,\,n\sigma^2)$, und folglich $\bar X \sim N(\mu,\,\sigma^2/n)$.},
    type=dinge
}

\newglossaryentry{formel_ki_bekannte_varianz}{
    name={Wie lautet das Konfidenzintervall für $\mu$ bei $n$ normalverteilten Stichprobenvariablen mit bekannter Varianz $\sigma^2$?},
    description={$\left[\bar X - \sigma\sqrt{\tfrac1n}\,c,\ \bar X+\sigma\sqrt{\tfrac1n}\,c\right]$ mit $c=z_{1-\alpha/2}$, sodass $P(\mu\in[\cdot])\ge 1-\alpha$.},
    type=dinge
}

%

\newglossaryentry{bem_intervallbreite_und_stichprobengroesse}{
    name={Wie verändert sich die Breite des Konfidenzintervalls mit wachsender Stichprobengröße $n$?},
    description={Das Intervall wird schmaler -- die Halbbreite $\sigma\sqrt{1/n}\,c$ sinkt mit wachsendem $n$.},
    type=dinge
}

\newglossaryentry{formel_versuchsplanung_stichprobengroesse}{
    name={Wie bestimmt man die minimale Stichprobengröße $n$, damit die Halbbreite des Konfidenzintervalls einen vorgegebenen Wert $d$ nicht überschreitet?},
    description={Man löst $\sigma\sqrt{1/n}\,c\le d$ nach $n$ auf und erhält $n \ge \left(\dfrac{\sigma c}{d}\right)^2$.},
    type=dinge
}

\newglossaryentry{def_konfidenzintervall_formal}{
    name={Wie ist ein Konfidenzintervall (Vertrauensintervall) formal definiert?},
    description={Seien $U,V$ Funktionen der Stichprobe $X_1,\dots,X_n$. $[U,V]$ heißt Konfidenzintervall für den unbekannten Parameter $\theta$ zum Konfidenzniveau $1-\alpha$, wenn $P_\theta(\theta\in[U,V])\ge 1-\alpha$ für alle $\theta$ gilt.},
    type=dinge
}

\newglossaryentry{bem_intuition_konfidenzniveau}{
    name={Was bedeutet das Konfidenzniveau $1-\alpha$ eines Konfidenzintervalls anschaulich?},
    description={Es ist die (garantierte Mindest-)Wahrscheinlichkeit, mit der das \emph{zufällige} Intervall $[U,V]$ den wahren, festen Parameter $\theta$ überdeckt -- die Zufälligkeit steckt im Intervall, nicht im Parameter.},
    type=dinge
}

\newglossaryentry{formel_ki_unbekannte_varianz}{
    name={Wie lautet das Konfidenzintervall für den Erwartungswert bei einer Stichprobe aus einer Normalverteilung mit unbekannter Varianz?},
    description={$\left[\bar X - \hat\sigma\sqrt{\tfrac1n}\,t_{n-1,1-\alpha/2},\ \bar X+\hat\sigma\sqrt{\tfrac1n}\,t_{n-1,1-\alpha/2}\right]$ mit $\hat\sigma^2=\tfrac{1}{n-1}\sum_i(x_i-\bar x)^2$.},
    type=dinge
}

%

\newglossaryentry{bem_warum_t_verteilung_statt_normal}{
    name={Warum genügt bei unbekannter Varianz nicht mehr das Quantil der Standardnormalverteilung, sondern das der Studentschen $t$-Verteilung?},
    description={Weil zusätzlich zum Erwartungswert auch die Varianz aus derselben Stichprobe geschätzt werden muss. Diese zusätzliche Unsicherheit durch die Varianzschätzung wird durch die (breiteren) Quantile der $t$-Verteilung berücksichtigt.},
    type=dinge
}

\newglossaryentry{def_studentsche_t_verteilung_ki}{
    name={Wie ist die Studentsche $t$-Verteilung mit $n-1$ Freiheitsgraden definiert?},
    description={Als die Verteilung der Zufallsvariable $T=\dfrac{\bar X-\mu}{\hat\sigma/\sqrt n}$.},
    type=dinge
}

\newglossaryentry{bem_eigenschaften_t_verteilung_ki}{
    name={Welche Eigenschaft hat die Studentsche $t$-Verteilung im Vergleich zur Standardnormalverteilung für wachsende Freiheitsgrade?},
    description={Für $n\to\infty$ (wachsende Freiheitsgrade) nähert sich die $t$-Verteilung der Standardnormalverteilung an.},
    type=dinge
}

\newglossaryentry{bem_bernoulli_varianzschaetzer_approximation}{
    name={Wie lässt sich bei einer Bernoulli-verteilten Stichprobe die (vom Parameter selbst abhängige) Varianz näherungsweise schätzen?},
    description={$\hat\sigma^2 \approx \bar X(1-\bar X)$ für große $n$.},
    type=dinge
}

\newglossaryentry{formel_ki_bernoulli}{
    name={Wie lautet das approximative Konfidenzintervall für die Erfolgswahrscheinlichkeit bei einer Bernoulli-verteilten Stichprobe?},
    description={$\left[\bar X - \sqrt{\tfrac{\hat\sigma^2}{n}}\,z_{1-\alpha/2},\ \bar X+\sqrt{\tfrac{\hat\sigma^2}{n}}\,z_{1-\alpha/2}\right]$, wobei für großes $n$ das $t$-Quantil durch das $z$-Quantil der Standardnormalverteilung ersetzt werden darf.},
    type=dinge
}

%

\newglossaryentry{bem_warum_z_statt_t_bei_bernoulli_grosses_n}{
    name={Warum darf man beim Konfidenzintervall für die Bernoulli-Erfolgswahrscheinlichkeit für großes $n$ das $t$-Quantil durch das $z$-Quantil ersetzen?},
    description={Weil sich die Studentsche $t$-Verteilung für wachsende Freiheitsgrade (also großes $n$) der Standardnormalverteilung annähert -- die Approximation wird für großes $n$ vernachlässigbar ungenau.},
    type=dinge
}
